\documentclass[10pt,a4paper]{amsart}
\usepackage[margin=19mm]{geometry}
\usepackage{amsmath,amssymb,mathtools}
\usepackage[scr=rsfs,cal=euler]{mathalfa}
\usepackage{xfrac}
\usepackage[unicode,hidelinks,bookmarks=false]{hyperref}
\newcommand{\ie}{i.e.\ }
\newcommand{\cf}{cf.\ }
\newcommand{\resp}{respectively\ }
\newcommand{\Subsection}{\subsection}
\providecommand{\nobreakdash}{\nobreak\hbox{-}}
\setlength{\emergencystretch}{2em}
\multlinegap0pt
\usepackage[shortlabels]{enumitem}
\usepackage{amssymb}
\newtheorem{theorem}{Theorem}
\newtheorem{definition}{Definition}
\title{Rigidity and index of free boundary minimal $Y$-cones in the unit ball}
\author{Elham Matinpour}
\date{Source: arXiv:2509.24137v2 (CC BY 4.0)}
\begin{document}
\maketitle

\noindent\emph{Source and licence.} Rigidity and index of free boundary minimal $Y$-cones in the unit ball. Version arXiv:2509.24137v2 (CC BY 4.0), distributed by arXiv under the Creative Commons Attribution 4.0 International licence, http://creativecommons.org/licenses/by/4.0/, as recorded in the arXiv licence metadata for this exact version and shown on the abstract page https://arxiv.org/abs/2509.24137v2. Full licence notice: \texttt{licenses/arxiv-2509.24137-CC-BY-4.0.txt}.\par\smallskip

\noindent\textit{Editorial scope.} Counted unit: the article's theorem thm:main2-n (source0.tex lines 158--163). The article's complete original proof of it is delivered with it (source0.tex lines 875--969, one \texttt{proof} environment of 95 lines, which the article places away from the statement). No mathematics is written, repaired or re-derived: the statement and the proof are the article's own lines, in the article's own notation, and no retained line is substituted or re-flowed.  Retained uncounted as the source-local context the proof needs: lines 143--146; lines 307--317.  Nothing was omitted from any retained span, so the package carries no omission record.  No retained line contains a citation, so the wrapper carries no bibliography and no bibliography entry.  No retained line contains a figure environment, an image-inclusion command or drawing code, so no image is delivered and the package declares no asset dependency.  Editorial formatting only, in addition: an AMS \texttt{amsart} preamble that restates the article's own declarations and macros used by the retained lines, namely the environments \texttt{theorem}, \texttt{definition} on separate counters, together with the standard packages those lines need.  The retained environments print in this document as Theorem 1, Definition 1 in order of appearance, whereas the article prints the same items with the numbering of its own full document.  The complete native source of the article as distributed in the arXiv e-print for this exact version is bundled unchanged as \texttt{sources/185926/source0.tex}.
\begin{theorem}\label{thm:main2}
The Morse index of the free boundary flat $Y$--cone in $\mathbb{B}^3$ 
is two and its nullity is five.
\end{theorem}

\begin{definition}[Flat $Y$–cone]\label{def:Ycone}
Let
\[
\hat D=\{(x,y)\in\mathbb{R}^2: x^2+y^2\le1,\ x\ge0\}
\]
be the half–disk. Let $\hat D_\theta$ denote the rotation of $\hat D$ by angle $\theta$ about the $y$–axis in $\mathbb{R}^n$.
The \emph{flat $Y$–cone} is defined by
\[
YC=\hat D\cup\hat D_{120}\cup\hat D_{-120}.
\]
\end{definition}

\begin{theorem}\label{thm:main2-n}
The Morse index of the free boundary flat $Y$--cone in $\mathbb{B}^n$ 
is $2(n-2)$ and its nullity is $5(n-2)$. Moreover, if 
$\Sigma \subset \mathbb{B}^n$ is a free boundary minimal $Y$--surface with 
Morse index $2(n-2)$, then $\Sigma$ must be the flat $Y$--cone.
\end{theorem}

\begin{proof}[Proof of Theorem \ref{thm:main2-n}]
We extend the proof of Theorem~\ref{thm:main2} to higher dimensions, highlighting the modifications needed for the general case.  \\ 
Let $YC \subset \mathbb{B}^n$ denote the flat $Y$-cone, obtained by gluing three planar half-disks along a common diameter at $120^\circ$ angles, as defined in \ref{def:Ycone}. Each face $\Sigma_j$ is flat, so the Jacobi operator on $\Sigma_j$ is
\[
J_j = \Delta_{\Sigma_j}, \quad |A_{\Sigma_j}|^2 = 0.
\]
The $J$-Steklov eigenvalue problem on each face depends only on the intrinsic two-dimensional geometry and the boundary conditions:
\begin{itemize}
    \item The eigenfunctions $w_k(r,\theta) = a_k r^k \cos(k\theta) + b_k r^k \sin(k\theta)$ have eigenvalues $\sigma_k = k$, independent of $n$.
    \item The compatibility condition along the junction, $f_1 + f_2 + f_3 = 0$ on $\Gamma$, is dimension-independent.
\end{itemize}

In the scalar (codimension one) case, the negative eigenspace of the index form
\[
Q(f,f) = \sum_{j=1}^3 \int_{\partial' \Sigma_j} (\sigma - 1) (f|_{\partial' \Sigma_j})^2
\]
is exactly two-dimensional, spanned by constant functions on the faces satisfying the compatibility condition.  

For $YC \subset \mathbb{R}^n$, the normal bundle has rank $n-2$. A direct computation using the explicit Steklov eigenfunctions shows that the two negative directions per normal component satisfy the junction compatibility independently and are linearly independent across the $n-2$ components. Therefore, the total Morse index is
\[
\operatorname{index}(YC) = 2(n-2).
\]

Now suppose $\Sigma \subset \mathbb{B}^n$ is a free-boundary minimal $Y$-surface with Morse index $2(n-2)$. Let $E_1,\dots,E_n$ denote the standard basis of $\mathbb{R}^n$, and let $NE_k$ be the normal projection of $E_k$ onto $\Sigma$. As in the three-dimensional case:
\begin{itemize}
    \item On each face $\Sigma_j$, $n^k_j = \langle E_k, N_j \rangle$ satisfies $J_j n^k_j = 0$.
    \item The index form satisfies $Q(NE_k,NE_k) < 0$ unless $NE_k \equiv 0$.
\end{itemize}
Since the Morse index is $2(n-2)$, at most $2(n-2)$ of the $NE_k$ can be linearly independent negative directions. Therefore, there exist $(n-2)$ independent parallel vector fields $E$ such that $NE \equiv 0$ on $\Sigma$, i.e., $(n-2)$ independent directions are everywhere tangent to $\Sigma$.  

This implies that $\Sigma$ is contained in a three-dimensional affine subspace $V \subset \mathbb{R}^n$.\\
Within $V$, $\Sigma$ is a two-dimensional free-boundary $Y$-surface in $\mathbb{B}^3$ with Morse index two. By Theorem~\ref{thm:main2}, the only such surface is the flat $Y$-cone. Hence $\Sigma$ coincides with the flat $Y$-cone in $V \subset \mathbb{R}^n$.
\medskip


Now, we compute the nullity is $5(n-2)$. We work facewise and componentwise in the normal bundle.  Let $\Sigma_1,\Sigma_2,\Sigma_3$ denote the three flat half-disk faces of $YC$.  Fix an orthonormal trivialization of the normal bundle on each face; because each face is flat and embedded as a planar half-disk, the Jacobi operator acts componentwise as the scalar Laplace operator $\Delta$ on each scalar normal component.  Thus a global normal Jacobi field may be written as a triple
\[
\Phi=(\phi_1,\phi_2,\phi_3),\qquad \phi_j:\Sigma_j\to\mathbb{R}^{n-2},
\]
with the understanding that $\phi_j$ is vector-valued (with $n-2$ scalar components) and each scalar component is harmonic on the interior of the corresponding half-disk.
\medskip

On each face $\Sigma_j$ the Jacobi equation $J\phi_j=0$ is the Laplace equation $\Delta \phi_j=0$ (componentwise). The linearized free boundary condition on the circular arc $\partial'\Sigma_j$ is exactly the Steklov condition
\[
\partial_\nu \phi_j = \phi_j \quad\text{on }\partial'\Sigma_j,
\]
where $\nu$ denotes the outward conormal of $\Sigma_j$ along the circular arc. Note that, this is the same boundary condition used in the index computation: the boundary term in the index form contains $(\sigma-1)\int_{\partial'\Sigma_j}(\phi_j)^2$ when a Steklov separation with eigenvalue $\sigma$ is used.  

Hence any scalar component of a global Jacobi field corresponds to a Steklov eigenfunction on the half-disk with eigenvalue $\sigma=1$.  It is classical, and follows by separation of variables or by homogeneity, that the Steklov eigenfunctions on the unit disk (and hence on the half-disk, restricted suitably) which have eigenvalue $\sigma=k$ are exactly the restrictions of homogeneous harmonic polynomials of degree $k$.  In particular, the $\sigma=1$ modes are precisely the degree-one harmonic polynomials, i.e., linear functions.  No higher degree ($k\ge2$) contributes to $\sigma=1$, and the constant functions correspond to $\sigma=0$.

Therefore every scalar component of a global Jacobi field on a face $\Sigma_j$ is necessarily the restriction of a linear function on that planar face.  Equivalently, in local coordinates $(s,t)$ on the half-disk where $s$ is the coordinate along the diameter $\Gamma$ and $t$ is the transverse coordinate (so $\Gamma=\{t=0\}$ and the circular arc is $s^2+t^2=1$, $t\ge0$), any scalar $\sigma=1$ solution on $\Sigma_j$ has the form
\[
f_j(s,t) = A_j s + B_j t,
\]
for constants $A_j,B_j\in\mathbb{R}$, depending on the chosen scalar component.
\medskip

Global admissibility across the three faces requires the usual junction compatibility condition along the common diameter $\Gamma$, namely
\begin{equation}\label{eq:compat-vector}
\phi_1|_{\Gamma} + \phi_2|_{\Gamma} + \phi_3|_{\Gamma} = 0
\end{equation}
as an equality of vector-valued functions along $\Gamma$.  Restricting the linear form $f_j(s,t)=A_j s + B_j t$ to $\Gamma$ (where $t\equiv0$) shows that the restriction depends only on the $A_j$ coefficient:
\[
f_j|_{\Gamma}(s) = A_j s.
\]
Hence the compatibility condition \eqref{eq:compat-vector} for a single scalar normal component reduces to the scalar identity
\[
(A_1 + A_2 + A_3)\, s \;=\; 0 \qquad\forall s\in[-1,1],
\]
and therefore
\[
A_1 + A_2 + A_3 = 0.
\]
When $\phi_j$ is vector-valued with $n-2$ scalar components, in the chosen normal frame, the compatibility condition is the vector equation obtained by applying the above scalar constraint componentwise; equivalently it imposes exactly $n-2$ independent scalar linear constraints (one per normal component) on the $3\times (n-2)$ collections of $A_j$-coefficients.
\medskip


For each face $j$ and for each scalar normal component there are two real parameters $(A_j,B_j)$ describing the degree-one harmonic solution on that face.  Thus for the three faces and a single scalar normal component there are $3\cdot 2=6$ parameters in total.  The compatibility condition along $\Gamma$ imposes one linear scalar constraint $A_1+A_2+A_3=0$, reducing the free parameters to $6-1=5$ for that scalar component.

Since the $n-2$ scalar normal components decouple (the Laplacian and the boundary condition act componentwise in the chosen trivialization), the same count applies independently to each normal component.  Therefore the full vector--valued kernel has dimension
\[
\underbrace{5}_{\text{degrees of freedom per scalar normal comp.}}\times \underbrace{n-2}_{\text{number of normal components}} \;=\; 5(n-2).
\]
\medskip


Finally, we must rule out any other possible kernel elements than those coming from the degree-one modes counted above.  Let $\Phi$ be any admissible global Jacobi field.  As observed in step (1), each scalar component of $\Phi$ is harmonic on the interior of each face and satisfies the Steklov condition $\partial_\nu \phi = \phi$ on the circular arcs; consequently each scalar component is a Steklov eigenfunction with eigenvalue $\sigma=1$ on each face.  By the classical homogeneous harmonic polynomial characterization of Steklov eigenfunctions, the only harmonic polynomials on the disk with Steklov eigenvalue $\sigma=1$ are linear polynomials.  Hence every scalar component is necessarily of the affine linear form $A_j s + B_j t$ on each face, and thus is included in the parameter family already counted.  No constants ($\sigma=0$) nor higher degree harmonics ($\sigma\ge2$) can contribute to the kernel.  This excludes any additional kernel elements.

\medskip

Combining the counts above yields $\dim\ker J=5(n-2)$.  That gives the general formula
\[
\dim\ker J = 5(n-2).
\]
\end{proof}

\end{document}
